% Figura: pipeline de limpieza por capas (sección 4).
\begin{tikzpicture}[
    font=\footnotesize,
    node distance=5mm and 20mm,
    caja/.style={draw, rounded corners=1pt, align=center, text width=21mm, minimum height=10mm, inner sep=3pt},
    etiqueta/.style={above, font=\scriptsize},
    flecha/.style={-Stealth, semithick}]
  \node[caja] (tlc) {TLC\\CDN oficial};
  \node[caja, right=of tlc] (bronze) {\textbf{Bronze}\\Parquet original};
  \node[caja, right=of bronze] (silver) {\textbf{Silver}\\Parquet limpio};
  \node[caja, right=of silver] (gold) {\textbf{Gold}\\CSV para Go};
  \node[caja, dashed, below=6mm of silver] (audit) {Auditoría\\SQL en DuckDB};
  \draw[flecha] (tlc) -- node[etiqueta] {descarga} (bronze);
  \draw[flecha] (bronze) -- node[etiqueta] {7 reglas} (silver);
  \draw[flecha] (silver) -- node[etiqueta] {features} (gold);
  \draw[flecha, dashed] (audit) -- (silver);
\end{tikzpicture}
